\documentclass{article}

\usepackage[UTF8]{ctex}
\usepackage{amsmath,amssymb}
\usepackage{geometry}
\usepackage{array}

\geometry{a4paper, margin=2.5cm}
\renewcommand{\arraystretch}{1.25}

\begin{document}

% ============================================================
\section*{强化学习基础主线整理：从决策建模到深度 Model-Free RL}

\subsection*{先讲一个故事}

学习强化学习时，最容易出现的误解是：把它看成一长串算法名字。

多臂老虎机、MDP、TD、DQN、Policy Gradient、Actor-Critic、PPO、SAC，听起来像一张越来越长的清单。初学者往往会问：「我是不是还差某个算法没学？」但真正困难的地方并不在于算法数量，而在于理解这些算法背后反复出现的几个问题。

\textbf{智能体如何描述一个决策问题？}

\textbf{如何估计一个动作的长期价值？}

\textbf{当状态空间太大时，如何用函数近似替代表格？}

\textbf{如果直接优化策略，如何控制梯度方差和更新幅度？}

\textbf{在连续控制中，如何同时保持探索、样本效率和训练稳定性？}

这组笔记目前覆盖的内容，正是在沿着这些问题逐步推进。它不是按算法流行度排列的，而是在构建一条从基础决策建模到现代深度 model-free RL 的理解路线。

\textbf{本文的目标}：按照概念功能回看这条学习路线，说明各章节之间的承接关系，并指出它们共同建立了怎样的强化学习理解框架。

% ============================================================
\section*{一条主线，而不是一组孤立算法}

强化学习的基本对象是一个会行动的智能体。它在环境中观察状态 $s$，选择动作 $a$，收到奖励 $r$，再进入下一个状态 $s'$。如果只看这一句话，RL 看起来并不复杂。但真正的问题在于：每个动作的好坏并不只由眼前奖励决定，而取决于它对未来整条轨迹的影响。

因此，强化学习的主线可以概括为：

\[
\boxed{
\text{描述决策问题}
\;\Longrightarrow\;
\text{估计长期价值}
\;\Longrightarrow\;
\text{稳定地优化策略}
}
\]

早期章节解决的是「如何把问题说清楚」。多臂老虎机把探索与利用的冲突放到最简单的单状态问题中；MDP 则把单状态扩展为状态转移，使状态、动作、奖励、转移概率和折扣回报成为统一语言。

中间章节解决的是「如何估计未来」。TD 学习引入 bootstrap，使智能体不必等到整局结束才能更新价值；n 步回报进一步展示 TD 和 Monte Carlo 之间不是二选一，而是一条可以调节的偏差--方差连续谱。

深度 RL 章节解决的是「如何让这些思想进入高维状态和连续动作」。DQN 用神经网络替代 Q 表，并通过经验回放和目标网络稳定训练；Policy Gradient 开启直接优化策略的路线；Actor-Critic 用价值函数降低策略梯度方差；PPO 用裁剪目标限制策略更新幅度；SAC 则在最大熵和 off-policy 框架下处理连续控制。

换句话说，这条路线并不是从一个算法跳到另一个算法，而是在持续回答同一个问题：智能体如何在不确定环境中，通过交互数据学习一个稳定、可改进的决策规则。

% ============================================================
\section*{当前内容地图}

下表按概念功能整理已有章节，而不是按算法名称堆叠。

\begin{center}
\begin{tabular}{p{3.2cm}|p{4.2cm}|p{7.0cm}}
\textbf{模块} & \textbf{对应章节} & \textbf{核心问题} \\
\hline
决策问题的语言 &
多臂老虎机，MDP &
什么是探索、利用、状态、动作、奖励和转移？ \\
\hline
长期价值估计 &
TD，n 步回报 &
如何在不等完整回合结束的情况下估计长期回报？ \\
\hline
深度价值学习 &
DQN &
当状态空间太大，Q 表装不下时，如何用神经网络近似价值函数？ \\
\hline
策略直接优化 &
Policy Gradient，Actor-Critic，PPO &
如何直接优化策略，并控制高方差和更新过大的问题？ \\
\hline
连续控制与最大熵 &
SAC &
在连续动作空间中，如何兼顾探索、样本效率和训练稳定性？
\end{tabular}
\end{center}

这张表显示，这些章节共同构成了一条相对完整的深度 model-free RL 学习路径。它从问题建模出发，经由价值估计进入函数近似，再过渡到策略优化和连续控制。

沿着这条路径，可以看到深度 RL 中最常见的两条路线：一条是以 Bellman 方程和价值函数为核心的价值学习路线，另一条是以期望回报梯度为核心的策略优化路线。Actor-Critic 则把二者结合起来：Actor 负责改进策略，Critic 负责提供价值估计和低方差信号。

% ============================================================
\section*{五个核心问题}

\subsection*{问题一：如何描述一个决策问题？}

多臂老虎机是强化学习中最小的决策模型。它没有状态转移，只有动作选择和奖励反馈。因此，它非常适合用来解释探索与利用：智能体既要利用当前看起来最好的动作，又要探索那些不确定但可能更优的动作。

MDP 则把问题推进到序列决策。一个动作不再只影响眼前奖励，还会改变未来状态分布。此时，智能体优化的不是单步奖励，而是折扣累计回报：

\[
G_t = r_t + \gamma r_{t+1} + \gamma^2 r_{t+2} + \cdots
\]

这一步非常关键。只有先把问题写成 MDP，后面的 Bellman 方程、价值函数、策略优化和深度 RL 才有共同的数学语言。

\subsection*{问题二：如何估计长期价值？}

如果智能体只看眼前奖励，它会变成短视的贪心系统；如果每次都等到整局结束再更新，又会带来很高方差。TD 学习的核心贡献，是用一步 bootstrap 把未来估计接到当前估计上：

\[
V(s_t) \leftarrow V(s_t) + \alpha\left[r_t + \gamma V(s_{t+1}) - V(s_t)\right]
\]

这里的 TD 误差

\[
\delta_t = r_t + \gamma V(s_{t+1}) - V(s_t)
\]

既是价值估计的修正信号，也会在 Actor-Critic 和 GAE 中反复出现。

n 步回报进一步说明，TD 和 Monte Carlo 并不是互相替代的两端，而是一条偏差--方差权衡曲线。$n$ 越小，更新越及时但偏差更大；$n$ 越大，目标更接近真实回报但方差更高。这个思想后来在 GAE 中以 $\lambda$ 的形式继续出现。

\subsection*{问题三：状态空间太大怎么办？}

Q-learning 在小型离散环境中可以维护一张 Q 表，但一旦状态变成连续向量或图像，表格方法就失效了。DQN 的关键转变，是用神经网络表示 $Q(s,a;\theta)$：

\[
Q(s,a;\theta) \approx Q^*(s,a)
\]

这一步把强化学习带入深度学习框架。但神经网络函数近似也带来训练不稳定：样本连续相关，目标值又会随着网络参数一起移动。DQN 的经验回放和目标网络，分别对应这两个问题。

经验回放把历史转换存入数据缓冲区，再随机采样训练，使样本分布更接近独立同分布；目标网络则让 Bellman 目标在短期内保持稳定，避免在线网络追逐一个每步都在移动的目标。

因此，DQN 的意义不只是「用神经网络玩游戏」，而是给深度价值学习提供了一套基本稳定化框架。

\subsection*{问题四：策略能否直接优化？}

价值函数路线先估计动作好坏，再通过 $\arg\max$ 选择动作。策略梯度路线则直接参数化策略 $\pi_\theta(a|s)$，并优化期望回报：

\[
J(\theta)=\mathbb{E}_{\tau\sim\pi_\theta}[G(\tau)]
\]

REINFORCE 用对数导数技巧得到可计算的梯度估计：

\[
\nabla_\theta J(\theta)
=
\mathbb{E}_{\tau\sim\pi_\theta}
\left[
G_t \nabla_\theta \log \pi_\theta(a_t|s_t)
\right]
\]

这条路线的优势是可以自然处理随机策略和连续动作，但问题是方差很高。Actor-Critic 的出现，正是为了解决这一点：Actor 负责更新策略，Critic 负责估计价值并提供低方差的优势信号。

PPO 在 Actor-Critic 的基础上进一步限制策略更新幅度。它不让新策略相对于旧策略变化过大，从而避免一次更新破坏已有策略。PPO 的核心并不是复杂，而是克制：每次只允许策略迈出有限的一步。

\subsection*{问题五：连续控制中如何保持稳定探索？}

在连续动作空间中，策略既要输出连续动作，又要保持足够探索。SAC 把这个问题放进最大熵框架：算法不只最大化奖励，也最大化策略熵。

\[
J(\pi) =
\sum_t
\mathbb{E}_{(s_t,a_t)\sim\rho_\pi}
\left[
r(s_t,a_t)
+ \alpha H(\pi(\cdot|s_t))
\right]
\]

这个目标使随机性不再只是训练早期的临时工具，而成为优化目标的一部分。SAC 同时结合 off-policy 数据复用、双 Q 网络、软目标更新和自动温度调节，使它成为连续控制任务中非常重要的基线算法。

到 SAC 为止，可以看到深度 RL 的一个核心模式：稳定训练往往不是靠单个公式完成的，而是由多个机制共同维持。价值估计、策略更新、样本复用、探索调节和目标平滑，缺一项都可能让训练变得脆弱。

% ============================================================
\section*{SAC 在主线中的位置}

SAC 很好地汇合了前面多条线索：价值函数估计、Actor-Critic 结构、off-policy 数据复用、目标网络平滑，以及探索与利用之间的权衡。

从价值函数角度看，DQN 展示了如何把 Bellman 更新、函数近似、经验回放和目标网络结合起来。从策略优化角度看，Policy Gradient、Actor-Critic 和 PPO 展示了如何从高方差梯度估计逐步走向稳定的策略更新。从连续控制角度看，SAC 进一步展示了最大熵、off-policy 学习和双 Q 估计如何组合成现代连续控制算法。

因此，当前内容已经回答了三个基础问题：

\begin{enumerate}
  \item 如何把强化学习问题写成统一的数学形式；
  \item 如何估计长期价值，并在偏差和方差之间取舍；
  \item 如何在深度神经网络条件下稳定训练价值函数和策略。
\end{enumerate}

在这个位置上回看前面的内容，可以得到一个比较清楚的结构：基础 model-free RL 的核心问题，已经从「如何估计价值」推进到「如何稳定地改进策略」，并进一步进入「如何在连续动作空间中保持有效探索」。

% ============================================================
\section*{总结}

这组强化学习笔记目前已经形成一条从基础决策建模到深度 model-free RL 的主线。多臂老虎机和 MDP 提供问题语言，TD 和 n 步回报提供价值估计框架，DQN 把价值学习推进到深度函数近似，Policy Gradient、Actor-Critic 和 PPO 建立策略优化路线，SAC 则把最大熵和 off-policy 学习带入连续控制。

回看这条路线，最重要的收获不是记住每个算法的名字，而是看到强化学习在反复处理几类核心矛盾：探索与利用，偏差与方差，价值估计与策略改进，样本效率与训练稳定性，以及随机性与最优性之间的关系。

理解这些矛盾，比单独记住某个算法的操作细节更重要。因为一旦抓住这些结构，再回到具体算法时，就能看清每个设计到底是在缓解哪一种不稳定、降低哪一种方差，或是在补偿哪一种估计偏差。

\subsection*{参考资料}

\begin{enumerate}
  \item Sutton, R.\,S.\ \& Barto, A.\,G.\ (2018). \textit{Reinforcement Learning: An Introduction} (2nd ed.). MIT Press.
  \item Mnih, V., et al.\ (2015). Human-level control through deep reinforcement learning. \textit{Nature}, 518, 529--533.
  \item Schulman, J., et al.\ (2017). Proximal Policy Optimization Algorithms. \textit{arXiv:1707.06347}.
  \item Haarnoja, T., et al.\ (2018). Soft Actor-Critic: Off-Policy Maximum Entropy Deep Reinforcement Learning with a Stochastic Actor. \textit{ICML}.
\end{enumerate}

\end{document}
